% Unnumbered research note. Numerical refinement only; not promoted or integrated.
\subsection*{Weighted final exceptional curves in necessary SNC refinement}

This note gives a proposed complete numerical transfer argument, subject to
independent actual-file checking.  It does not construct an order-module
atlas, compare root gerbes, or control thick end overlaps.  All blowups
below are at non-SNC points of the current reduced total boundary.  No
additional blowup at a smooth boundary point is permitted.

\paragraph{Local statement.}
Let a reduced curve germ in a complete regular equicharacteristic surface
over an algebraically closed field have analytic branches $f_1,\ldots,f_s$.
Put $\delta_i=\delta(f_i)$ and $I_{ij}=I(f_i,f_j)$, and give these
branches positive integral coefficients $c_i$.  Resolve the reduced total
boundary by the stipulated point blowups, stopping when it is SNC.
For a new exceptional curve $E$, put
$c_E=\sum_i c_i v_E(f_i)$.  Then the final new curves of square $-1$
satisfy
\[
 \sum_{E^2=-1,\ E\ \mathrm{new}}c_E
 \ \le\ 6\sum_i c_i\delta_i
       +2\sum_{i<j}(c_i+c_j)I_{ij}.                 \tag{D.WS.1}
\]
The assertion concerns the final $(-1)$-curves, not the sum of the
coefficients on all intermediate exceptional curves.  The latter can
be quadratic in a tangency order.

\paragraph{Two branch identities in arbitrary characteristic.}
Write $a_E$ for the coefficient of $E$ in the relative canonical divisor
over the original regular surface.  Suppose the final strict transform
of branch $i$ meets $E$ transversely, at a point on no other exceptional
curve or strict branch.  Then
\[
 v_E(f_i)-a_E=2\delta_i,\qquad
 v_E(f_j)=I_{ij}\quad(j\ne i).                       \tag{D.WS.2}
\]
Here is the precise delta/adjunction calculation for the first identity.
At a blowup where the current strict branch has multiplicity $n$, its
strict transform is $D'=\pi^*D-nF$.  The ideal sequence and
$\pi_*\mathcal O(nF)=\mathcal O$,
$\operatorname{length}R^1\pi_*\mathcal O(nF)=n(n-1)/2$
give
\[
 \delta(D)-\delta(D')=n(n-1)/2.
\]
The curve map here is finite birational: its exceptional fibre is finite,
so the proper strict-transform map is quasi-finite.  Normalization is
unchanged.  Adjunction on the two regular surfaces gives
\[
 \omega_{D'}=\pi^*\omega_D\otimes
                         \mathcal O((1-n)F)|_{D'}.
\]
Pulling this identity to the eventual normalized branch, the equation
of $F$ has order $n$: in a blowup chart it is the parameter of minimum
order on the current branch.  Thus this step contributes $n(n-1)$ to
the discrepancy between the pulled-back branch dualizing line and its
normalization dualizing line.  Sum along the branch.  At the final
transverse point the same discrepancy, computed directly from
$K_{Y'/Y}=\sum a_EE$ and
$\pi^*D_i=D_i'+\sum v_E(f_i)E$, is $v_E(f_i)-a_E$.
The delta drops sum to $\delta_i$, proving the first identity.
This uses adjunction, not a derivative or a separable Puiseux expansion.
The second identity follows by restricting $f_j$ to the normalized
branch $i$: at its final point every other total-transform component
is a unit, and an equation of $E$ has order one.

The quadratic-transformation facts just used are the complete written
Statements and proofs of Stacks Tags 0AGQ, 0AGR and 0AGS, Lemmas
54.3.1, 54.3.2 and 54.3.4.  Their displayed cohomology calculation is
valid in every characteristic.  The explicit omitted verification in
0AGS concerning the graded exceptional-fibre map is retained as a
Known input.  Cartier adjunction and the canonical formula for a point
blowup are the same known inputs used in the previously audited cycle
and SNC-refinement arguments.  No new conductor theorem or
characteristic-zero resolution theorem is being silently invoked.

\paragraph{The intrinsic single-branch discrepancy bound.}
First resolve one singular branch alone, stopping once it is smooth and
transverse to exactly one exceptional component, away from exceptional
nodes.  Write $m_E=v_E(f_i)$, $d_E=m_E-a_E$ and
$q_E=2d_E-a_E$.  At a centre of strict-branch multiplicity $n$, whose
old exceptional parents form a set $J$ of size at most two,
\[
 a_{\rm new}=1+\sum_{E\in J}a_E,\quad
 d_{\rm new}=n-1+\sum_{E\in J}d_E,\quad
 q_{\rm new}=2(n-1)-1+\sum_{E\in J}q_E.              \tag{D.WS.3}
\]
While the branch is singular, $n\ge2$, so induction gives $q_E\ge1$
for all exceptional components then created.  The branch never becomes
singular again after it first becomes smooth.  The remaining cases are
as follows.

If the first smooth centre has one exceptional parent, it is a tangency;
otherwise the single branch is already resolved there.  Its parent has
$q\ge1$, so the first smooth blowup gives $q_{\rm new}\ge0$.
Thereafter the branch is at the corner between that fixed tangent
parent and the newest exceptional curve.  It is transverse to the
newest curve.  If still tangent, its tangent direction is the fixed
old parent.  Each further corner blowup has
$q_{\rm new}=q_{\rm old}+q_{\rm recent}-1\ge q_{\rm recent}$.
This continues until the last corner blowup makes the branch meet the
new exceptional curve away from its old parents.

If the first smooth centre has two exceptional parents, both have
$q\ge1$.  The first new margin is therefore at least one.  If the
branch is transverse to both, that blowup resolves the corner.  If
it is tangent to one, retain that parent and apply the preceding
corner argument.  A smooth branch cannot be tangent to both transverse
parents.  These are all cases, since the exceptional divisor alone is
SNC throughout.  They can also be read directly from the local
equation $x=y^e\cdot\text{unit}$ for a smooth branch tangent to
$x=0$: each blowup reduces $e$ by one, and the last corner is removed
by one final blowup.  This argument does not divide by $e$.

Consequently the last exceptional curve met in the intrinsic resolution
has $q_E\ge0$.  By (D.WS.2), $d_E=2\delta_i$, and hence
\[
             a_E\le4\delta_i,\qquad m_E\le6\delta_i.
                                                               \tag{D.WS.4}
\]
For an initially smooth branch take the empty intrinsic resolution and
initial values $a=m=\delta_i=0$.

\paragraph{Other branches and the final leaves.}
The centres on branch $i$ before it is intrinsically resolved are
precisely the centres forced by that branch and the exceptional
divisor.  Blowups at other points do not change this sequence locally.
Once the branch is intrinsically resolved, a further necessary centre
on it must contain another strict branch.  Such a centre has exactly
one exceptional parent on branch $i$.  After blowing up, branch $i$
is again smooth and transverse to the new exceptional, away from the
old parent.  Thus all these extra centres are free along this branch,
and each adds one to both $a_E$ and $m_E$ along its path.

Each extra centre consumes at least one unit of intersection with
another strict branch: the intersection drop is the product of the
two strict multiplicities, and branch $i$ now has multiplicity one.
The number $L_i$ of extra centres therefore obeys
\[
 L_i\le\sum_{j\ne i}I_{ij},\qquad
 v_E(f_i)\le6\delta_i+L_i.                           \tag{D.WS.5}
\]
This includes an initially smooth branch: the first blowup involving it
is an extra centre, with no old exceptional parent, and creates $a=m=1$.
Every subsequent extra centre has the one-parent form just described.

A final new $(-1)$-curve was never blown up after its creation.
Its creating centre contained a strict original branch, because the
exceptional divisor alone is SNC.  At least one such branch therefore
still meets it.  At the final SNC stage that branch meets it transversely
away from all other boundary components.  Different final $(-1)$-curves
can be assigned different analytic branches, since each normalized
branch over the original point has only one point on the exceptional
fibre.  For a curve assigned to $i$, (D.WS.2) and (D.WS.5) give
\[
 c_E\le6c_i\delta_i+\sum_{j\ne i}(c_i+c_j)I_{ij}.
\]
Sum over the distinct assigned branches, and enlarge to all branches.
Each unordered pair is counted at most twice.  This proves (D.WS.1).
An ordinary self-node of a global irreducible component is treated as
two distinct analytic branches; the one necessary blowup for strict
normal crossings is included by exactly the same argument.

\paragraph{Charging only the original bad points.}
Now take the ORIGINAL relative-minimal cover model with anti-nef
$C=\sum c_iD_i$ and the core of the existing degree-linear core and
weighted-chain propositions.  All noncore curves are smooth rational
$(-2)$-curves.  The already proved estimates, with constants depending
only on the fixed base, give
\[
 \sum_{i\in\mathcal K}c_i(b_i+p_a(D_i)+1)=O(t),
 \quad
 \sum_{i\in\mathcal K}c_i\operatorname{val}_i=O(t),
 \quad b_i=-D_i^2.                                  \tag{D.WS.6}
\]
The valence counts intersection multiplicity.  The second assertion
follows rowwise from the same Schur estimate in the core proof:
the noncore attachment count at $i$ is less than $2b_i$; weighted
core--core incidences follow by summing anti-nef inequalities.  Also
\[
 \sum_{i\in\mathcal K}\sum_{j\ne i}c_j(D_i\cdot D_j)
                  \le\sum_{i\in\mathcal K}c_i b_i=O(t).
                                                               \tag{D.WS.7}
\]
Apply the local bound only at the original non-SNC points.  Intrinsic
branch deltas and intersections of analytic branches of one global
component are bounded by its normalization defect, at most its
arithmetic genus; every singular global component is in the core.
Pairs of distinct global components with at least one core component
are paid by (D.WS.6)--(D.WS.7).

There cannot be three distinct noncore curves at one point: their
principal negative-intersection matrix has diagonal two and every
off-diagonal at most minus one, contradicting positive definiteness
on the all-one vector.  Nor can two noncore curves have intersection
multiplicity at least two, by the corresponding two-by-two principal
minor.  Hence a bad point containing a pair of noncore curves must
also contain a core branch.  Charge the two weights of that pair to
their intersections with that core branch.  There are at most two
noncore branches at the point, so this does not cause unbounded
overcounting.  A point containing only two transverse noncore curves
is already SNC and is not included.  In particular the many good
nodes of a long $(-2)$-chain are never charged.

The identity for an integral curve's delta is the sum of its analytic
branch deltas plus the pairwise intersections of its distinct analytic
branches.  Thus the same-component pair terms in (D.WS.1) can be
absorbed into $6c_i\delta(D_i)$ (their coefficient is $4c_i$).
These observations show that the sum of the local right sides is
$O(t)$.  Therefore on the necessary SNC refinement,
\[
       \sum_{E\ \mathrm{new},\ E^2=-1}c_E=O(t).        \tag{D.WS.8}
\]

\paragraph{The refined numerical core, including zero-$Z$ $(-1)$-curves.}
Let $C'=\rho^*C$ and $Z'=\rho^*Z$, so $C'\le Z'\le\alpha C'$.
The pullback identities already proved give
$K_{Y'}\cdot C'=K_Y\cdot C=O(t)$, with a two-sided bound.
On the SNC model put $\kappa_i=K_{Y'}\cdot D_i=2g_i-2+b_i$.
Its only negative values are $-1$ on smooth rational $(-1)$-curves.
Original point-mapped $(-1)$-curves were contracted; original radial
ones have total weighted mass $O(t)$ by the old core budget.  New
ones have the bound (D.WS.8).  It follows from the canonical
intersection identity that
\[
                 \sum_{\kappa_i>0}c_i\kappa_i=O(t).
\]
Define the refined core to contain all radial strict transforms, all
vertices with $\kappa_i>0$, and all new $(-1)$-curves.  In particular
the latter are retained even though $Z'\cdot E=0$.  Adjunction and
the radial weight bound give
\[
           \sum_{i\in\mathcal K'}c_i(g_i+b_i+1)=O(t).
\]
Every remaining curve is a smooth rational $(-2)$-curve.  Apply the
same rowwise Schur estimate and the anti-nef inequalities for $C'$.
The result is the desired numerical transfer
\[
 \sum_{i\in\mathcal K'}Z'_i
          (g_i+b_i+\operatorname{val}_i+1)=O(t).      \tag{D.WS.9}
\]
This does not control products of adjacent cycle coefficients, replace
the new $(-1)$-curves by ADE blocks, or prove that all remaining
components are two-ended chains.  Forks and core attachments remain.
It supplies no order-valued vanishing, small projective presentation,
relative saturation, or unbounded-period representation count.

\paragraph{Read scope and status.}
The complete current proofs of the degree-linear core, weighted-chain
compression and controlled SNC refinement were reread.  The full
written Stacks 0AGQ/0AGR/0AGS proofs were freshly opened and read on
20 November 2026; their explicitly omitted routine details remain
known inputs.  The branch recurrences, margin cases and weighted
charging argument are supplied here for independent checking.
This note is unnumbered and unpromoted.  Even if every numerical
step passes, the main hereditary order-module atlas gate stays open.
